\documentclass[tikz,border=3mm,multi=tikzpicture]{standalone}
\usepackage{eleve-fisica}

\begin{document}

% FIGURA 1 — fig-01-equilibrio-estatico-e-dinamico
\begin{tikzpicture}[fisica figura, x=1cm, y=1cm]
  \path[use as bounding box] (-0.2,-0.2) rectangle (10.8,10.0);
  \node[font=\Large\bfseries, text=eleveazul] at (5.15,9.45)
    {equilíbrio estático};
  \coordinate (A) at (5.15,7.0);
  \node[fisica corpo, minimum width=1.65cm, minimum height=1.1cm] at (A) {};
  \draw[fisica forca] (A) -- ++(0,1.65)
    node[fisica rotulo, above] {$\vec N$};
  \draw[fisica forca] (A) -- ++(0,-1.65)
    node[fisica rotulo, below] {$\vec P$};
  \node[fisica medida] at (2.0,7.0) {$\vec v=0$};

  \node[font=\Large\bfseries, text=eleveazul] at (5.15,4.4)
    {equilíbrio dinâmico};
  \coordinate (B) at (5.15,1.95);
  \node[fisica corpo, minimum width=1.65cm, minimum height=1.1cm] at (B) {};
  \draw[fisica forca] (B) -- ++(0,1.65)
    node[fisica rotulo, above] {$\vec N$};
  \draw[fisica forca] (B) -- ++(0,-1.65)
    node[fisica rotulo, below] {$\vec P$};
  \draw[fisica movimento] ($(B.north)+(0,0.7)$) -- ++(2.2,0)
    node[fisica rotulo, right] {$\vec v$ constante};
  \node[font=\large\bfseries, text=elevelaranja] at (2.1,3.75)
    {$\vec F_R=0$ nos dois casos};
\end{tikzpicture}

% FIGURA 2 — fig-02-inercia-do-passageiro
\begin{tikzpicture}[fisica figura, x=1cm, y=1cm]
  \path[use as bounding box] (-0.2,-0.2) rectangle (11.2,9.4);
  \draw[fisica superficie] (0.55,1.05) -- (10.55,1.05);
  \draw[eleveazul, fill=eleveazulclaro, line width=1.4pt, rounded corners=5pt]
    (1.15,2.0) rectangle (9.3,5.25);
  \draw[elevecinza, line width=1.4pt] (2.3,2.0) -- (2.3,3.6);
  \node[fisica corpo, circle, minimum size=0.7cm] (P) at (5.5,3.65) {};
  \draw[elevecinza, line width=1.4pt] (4.6,2.35) -- (5.2,3.45);
  \draw[elevelaranja, line width=1.7pt] (4.6,4.75) -- (5.35,3.8);
  \draw[elevecinza, fill=white, line width=1.5pt] (2.7,1.05) circle (0.65);
  \draw[elevecinza, fill=white, line width=1.5pt] (7.9,1.05) circle (0.65);
  \draw[fisica forca] (P.center) -- ++(-2.1,0);
  \node[fisica rotulo] at (3.5,2.75) {$\vec F_{\text{cinto}}$};
  \draw[fisica movimento] (6.3,6.2) -- ++(2.6,0);
  \node[fisica rotulo, anchor=east] at (5.95,6.2)
    {tendência de manter $\vec v$};
  \draw[fisica movimento] (8.9,7.45) -- ++(-2.6,0);
  \node[fisica rotulo, anchor=east] at (5.95,7.45) {carro freia};
  \node[font=\large\bfseries, text=elevecinza] at (5.35,8.75)
    {o cinto fornece a força que reduz a velocidade};
\end{tikzpicture}

% FIGURA 3 — fig-03-resultante-de-forcas-perpendiculares
\begin{tikzpicture}[fisica figura, x=1cm, y=1cm]
  \path[use as bounding box] (-0.2,-0.2) rectangle (8.6,8.8);
  \coordinate (O) at (1.6,1.3);
  \draw[fisica forca] (O) -- ++(4.2,0);
  \node[fisica rotulo, below=8pt] at ($(O)+(2.1,0)$) {$F_x=30\,\mathrm{N}$};
  \draw[fisica forca] ($(O)+(4.2,0)$) -- ++(0,5.6);
  \node[fisica rotulo, anchor=west] at ($(O)+(4.55,2.8)$)
    {$F_y=40\,\mathrm{N}$};
  \draw[fisica forca] (O) -- ++(4.2,5.6);
  \node[fisica rotulo] at (2.42,5.06) {$F_R=50\,\mathrm{N}$};
  \draw[elevelaranja, line width=1.1pt] ($(O)+(3.65,0)$)
    -- ++(0,0.55) -- ++(0.55,0);
  \node[font=\large\bfseries, text=elevecinza] at (3.7,8.1)
    {triângulo $3$--$4$--$5$};
\end{tikzpicture}

% FIGURA 4 — fig-04-velocidade-e-aceleracao-opostas
\begin{tikzpicture}[fisica figura, x=1cm, y=1cm]
  \path[use as bounding box] (-0.2,-0.2) rectangle (10.4,7.9);
  \draw[fisica superficie] (0.65,1.1) -- (9.75,1.1);
  \node[fisica corpo, minimum width=1.8cm, minimum height=1.05cm] (A)
    at (5.0,2.0) {};
  \draw[fisica movimento] ($(A.north)+(0,1.0)$) -- ++(3.0,0)
    node[fisica rotulo, right] {$\vec v$};
  \draw[fisica forca] ($(A.north)+(0,2.3)$) -- ++(-2.5,0)
    node[fisica rotulo, left] {$\vec F_R$};
  \draw[fisica movimento] (5.0,5.9) -- ++(-2.0,0)
    node[fisica rotulo, left] {$\vec a$};
  \node[font=\Large\bfseries, text=eleveazul] at (5.0,7.25)
    {move para a direita e perde rapidez};
\end{tikzpicture}

% FIGURA 5 — fig-05-acao-e-reacao-em-corpos-diferentes
\begin{tikzpicture}[fisica figura, x=1cm, y=1cm]
  \path[use as bounding box] (-0.2,-0.65) rectangle (11.2,8.4);
  \draw[fisica superficie] (0.65,1.1) -- (10.55,1.1);
  \node[fisica corpo, minimum width=2.0cm, minimum height=2.5cm] (P)
    at (2.25,2.4) {pessoa};
  \node[fisica corpo, minimum width=2.5cm, minimum height=1.55cm] (C)
    at (8.25,1.9) {carrinho};
  \draw[fisica forca] (P.east) -- ++(2.4,0)
    node[fisica rotulo, midway, above=8pt] {$\vec F_{P\to C}$};
  \draw[fisica forca] (C.west) -- ++(-2.4,0);
  \node[fisica rotulo] at (5.25,0.35) {$\vec F_{C\to P}$};
  \node[fisica rotulo] at (2.25,5.45) {força sobre a pessoa};
  \node[fisica rotulo] at (8.25,5.45) {força sobre o carrinho};
  \draw[fisica auxiliar, ->] (2.25,5.0) -- (3.2,3.2);
  \draw[fisica auxiliar, ->] (8.25,5.0) -- (7.25,2.95);
  \node[font=\Large\bfseries, text=elevelaranja] at (5.25,7.45)
    {mesmo módulo; corpos diferentes};
\end{tikzpicture}


% FIGURA 6 — fig-06-diagrama-de-corpo-livre
\begin{tikzpicture}[fisica figura, x=1cm, y=1cm]
  \path[use as bounding box] (1.0,1.7) rectangle (9.4,11.4);
  \node[font=\Large\bfseries, text=eleveazul] at (5.2,11.0) {cena física};
  \draw[fisica superficie] (1.6,8.0) -- (9.0,8.0);
  \node[fisica corpo, minimum width=1.8cm, minimum height=1.15cm] (K)
    at (4.6,8.575) {};
  \draw[fisica forca] (K.center) -- ++(2.3,1.3);
  \node[fisica rotulo, anchor=west] at (7.1,9.95) {$\vec F$};

  \node[font=\Large\bfseries, text=eleveazul] at (5.2,6.9)
    {diagrama de corpo livre};
  \coordinate (O) at (4.6,4.1);
  \node[fisica ponto] at (O) {};
  \draw[fisica forca] (O) -- ++(0,1.3)
    node[fisica rotulo, above] {$\vec N$};
  \draw[fisica forca] (O) -- ++(0,-1.6)
    node[fisica rotulo, below] {$\vec P$};
  \draw[fisica forca] (O) -- ++(2.08,1.2);
  \node[fisica rotulo, anchor=south west] at (6.95,5.45) {$\vec F$};
  \draw[fisica componente] (O) -- ++(2.08,0);
  \node[fisica rotulo, anchor=north] at (5.64,3.80) {$F_x$};
  \draw[fisica componente] ($(O)+(2.08,0)$) -- ++(0,1.2);
  \node[fisica rotulo, anchor=west] at (6.95,4.7) {$F_y$};
  \draw[fisica eixo] (1.5,2.4) -- ++(1.2,0)
    node[fisica rotulo, right] {$x$};
  \draw[fisica eixo] (1.5,2.4) -- ++(0,1.2)
    node[fisica rotulo, above] {$y$};
\end{tikzpicture}


% FIGURA 7 — fig-07-par-livro-e-terra
\begin{tikzpicture}[fisica figura, x=1cm, y=1cm]
  \path[use as bounding box] (1.3,0.3) rectangle (6.7,6.0);
  \node[fisica corpo, minimum width=2.0cm, minimum height=0.9cm] (L)
    at (4.2,5.2) {};
  \node[fisica rotulo, anchor=east] at (2.95,5.2) {livro};
  \draw[fisica forca] (L.center) -- ++(0,-1.8);
  \node[fisica rotulo, anchor=east] at (3.95,4.075) {$\vec F_{T\to L}$};
  \begin{scope}
    \clip (1.4,0.4) rectangle (7.0,1.6);
    \draw[eleveazul, fill=eleveazulclaro, line width=1.6pt]
      (4.2,-2.5) circle (3.6);
  \end{scope}
  \node[fisica rotulo, anchor=east] at (2.5,0.85) {Terra};
  \draw[fisica forca] (4.2,1.1) -- ++(0,1.8);
  \node[fisica rotulo, anchor=west] at (4.45,2.0) {$\vec F_{L\to T}$};
\end{tikzpicture}

\end{document}
